\section{来源与核心问题：好模型的幕后不是一条训练命令}

旧稿没有填写 canonical video URL，只把 YouTube thumbnail 重复当成唯一课程图，并把 71 页官方 deck 压成短篇概述。本讲现以讲者个人官网直链的 Google Slides、Stanford Online 官方录像 \texttt{jm2hyJLFfN8}、官方 \texttt{en-US} 人工字幕和 BigCode/StarCoder primary sources 为依据。全部 71 页 slides 已落盘；其中 58 页包含独立教学内容，其余仅为分隔、动画中间态、完成标记、重复页或致谢。

\subsection{课程不是“数据最重要”一句话，而是问题链}

标题把 StarCoder 标成 use case，意味着课程的目标不是复述一个模型配置，而是用代码模型回答一般预训练问题：要多少数据、从哪里找、如何过滤、如何治理、怎样评估、何时相信 benchmark。后文每一章都沿着这个 question--answer 结构展开。

\lecturefigure{slide-01.jpg}{Behind the Scenes of LLM Pretraining：StarCoder use case}{Loubna Ben Allal 官方 deck，第 1 页}

\lecturefigure{slide-02.jpg}{课程路线：用一连串问题回答“怎样训练好 LLM”}{Loubna Ben Allal 官方 deck，第 2 页}

\begin{importantbox}{本讲的六个交付问题}
\begin{enumerate}
  \item 在给定 compute 下，模型参数与训练 token 怎样分配？
  \item raw web/code/synthetic sources 怎样转成可训练 corpus？
  \item filter 直觉怎样通过小模型 ablation 验证？
  \item license、opt-out、PII 与 decontamination 如何进入 pipeline？
  \item repository/file/commit/notebook 怎样格式化成训练 sequence？
  \item leaderboard 如何避免 contamination、单一 benchmark 与时间泄漏？
\end{enumerate}
\end{importantbox}

\teachervoice{00:00:50--00:02:10，Ben Allal 说整场报告只有一个“很 loaded”的问题，slides 会不断提出子问题再回答。StarCoder 是把抽象原则落到真实项目的案例。}

\subsection{Open-model 进展：权重开放带来哪些实际能力}

2023 年流传的 ``No Moat'' memo 代表一种判断：闭源领先并不意味着开源永远无法追赶。课堂引用它不是为了宣布竞争结束，而是解释为什么公开 weights 会产生新的 deployment/fine-tuning ecosystem。

\lecturefigure{slide-03.jpg}{“No moat”判断：开源社区逐步补齐训练拼图}{Loubna Ben Allal 官方 deck，第 3 页}

开放 weights 后，quantization 可以降低 memory footprint，本地 runtime 可以在 consumer hardware 上部署，parameter-efficient fine-tuning 可以针对组织数据做定制。这些价值不等于模型在所有任务上达到闭源系统能力，而是\textbf{可检查、可改造、可私有部署的选择空间更大}。

\lecturefigure{slide-04.jpg}{2024 年课堂示例：接近 GPT-4 评价的模型可在大型消费级桌面运行}{Loubna Ben Allal 官方 deck，第 4 页}

\begin{warningbox}{“性能接近”必须绑定版本、指标与日期}
slide 是 2024 年 5 月快照。Arena、MMLU、coding、安全、长上下文、tool use、latency 与 cost 会给出不同排序；本讲不把某张表上的接近改写成通用 parity。
\end{warningbox}

\teachervoice{00:02:10--00:03:20，讲者强调 open weights 的直接收益是 quantization、consumer deployment 与社区 fine-tuning，而不是一句抽象的“开源更好”。}

\subsection{Open releases 在增加，但训练透明度仍不足}

Gemma、Mistral 等公开模型让 open ecosystem 扩大；Arena 图也显示 2023 到 2024 年 open/closed gap 缩小。读这两页时应先看它们回答的问题：第一页是供给增长，第二页是特定人类偏好评测协议下的趋势。

\lecturefigure{slide-05.jpg}{开放模型发布增加：Gemma、Mistral 等进入竞争}{Loubna Ben Allal 官方 deck，第 5 页}

\lecturefigure{slide-06.jpg}{May 2024 Arena 视角：open/closed 模型差距缩小}{Loubna Ben Allal 官方 deck，第 6 页}

\begin{knowledgebox}{Arena 证据能说明什么}
Arena 用用户 pairwise preference 形成 Elo-like ranking，适合衡量交互式回答偏好；它不能单独衡量训练数据合法性、事实可靠性、rare-language coverage、code execution correctness 或部署成本。趋势可参考，能力边界仍需多维评估。
\end{knowledgebox}

\subsection{为什么公开权重却不公开数据与 recipe}

讲者给出两个现实原因：一是 disclosure 可能引发 license/copyright scrutiny，二是团队希望保留 competitive edge。这使模型 weights 可以下载，却无法回答样本从何而来、filter 怎么选、token 数怎样配、evaluation 是否泄漏。

\lecturefigure{slide-07.jpg}{公开模型的限制：数据与训练细节经常缺失}{Loubna Ben Allal 官方 deck，第 7 页}

\begin{warningbox}{不可复现不是“文档少一点”}
若缺少 dataset version、dedup threshold、sampling weights、tokenizer、training tokens、optimizer schedule、benchmark decontamination 和 seeds，研究者无法定位性能来自 architecture、data 还是评测泄漏，也无法可靠复用治理选择。
\end{warningbox}

\teachervoice{00:04:10--00:05:05，讲者把不披露的法律与竞争原因直接说出。她并未因此放弃总结公共规律，而是强调需要从多个公开 release 拼出 evidence。}

\subsection{Model、GPUs、Data：三角形里本讲选择 data}

Transformer、Mamba、MoE 都可能成为 model 选项，GPU 决定 budget 上限；但在相近 architecture 与 training technique 下，数据质量和 mixture 常成为固定预算内的主要差异。这里的“data backbone”是课程的 focus，不是声称 architecture 和 hardware 无关。

\lecturefigure{slide-11.jpg}{训练好 LLM 的三类资源：Model、GPUs 与 Data}{Loubna Ben Allal 官方 deck，第 11 页}

\begin{knowledgebox}{把三者写成约束优化}
可把训练目标粗略写为
\[
\min_{\theta,\mathcal{D},r}\ L(\theta;\mathcal{D},r)
\quad\text{s.t.}\quad
C_{\text{train}}\le B,\ C_{\text{serve}}\le S,
\]
其中 $\mathcal{D}$ 是数据与 mixture，$r$ 是训练 recipe，$B,S$ 分别是训练与服务预算。课程主要优化 $\mathcal{D}$，但约束仍来自模型和 GPU。
\end{knowledgebox}

\teachervoice{00:05:05--00:06:20，讲者说 architecture/techniques 越来越相似，所以“给定预算下”数据让模型拉开差距。限定词是 given budget，而不是任何规模都只看数据。}

\subsection{本章小结}

Open weights 扩大部署和定制空间，但训练透明度仍是主要缺口。课程因此把预训练还原成受 compute、serving、data、governance 和 evaluation 共同约束的工程，而 StarCoder 提供可检查的端到端案例。

